% !TeX program = xelatex
% UTF-8; faithful transcription, not a correction of the medical source.
% Bold text preserves yellow highlighting. Uncertain readings are explicit.
% Original photographs, in handwritten page order:
% pic01.jpg = 20261012_193712426.jpg
% pic02.jpg = 20261012_193720146.jpg
% pic03.jpg = 20261012_193725028.jpg
% pic04.jpg = 20261012_193730641.jpg
% pic05.jpg = 20261012_193735445.jpg
% pic06.jpg = 20261012_193741944.jpg
% pic07.jpg = 20261012_193748496.jpg
% pic08.jpg = 20261012_193754605.jpg
% pic09.jpg = 20261012_193758596.jpg
\newif\ifPNDStandalone
\ifdefined\settextfont
  \PNDStandalonefalse
\else
  \PNDStandalonetrue
\fi
\ifPNDStandalone
  \documentclass[a4paper,12pt]{extbook}
  \usepackage[margin=2.5cm]{geometry}
  \usepackage{amsmath,amssymb}
  \usepackage{enumitem}
  \usepackage{xepersian}
  \settextfont{Amiri}
  \setlatintextfont{FreeSerif}
\fi
\ifPNDStandalone
  \begin{document}
\fi
\begingroup
\emergencystretch=3em
\newcommand{\PNDen}[1]{\lr{#1}}
\newcommand{\PNDhl}[1]{\textbf{#1}}
\newcommand{\PNDunclear}[1]{[خوانش نامطمئن: #1]}
\newcommand{\PNDpage}[1]{%
  \clearpage\section*{صفحهٔ #1}
  \noindent\PNDen{Date: / / \quad Sat. Sun. Mon. Tue. Thu. Wed. Fri.}\par
  \noindent\PNDen{Subject:} ۱۵ ـ بیماری‌های اعصاب محیطی\par\medskip}
\chapter{بیماری‌های اعصاب محیطی}

\PNDpage{۱}
\section*{علائم بالینی نوروپاتی}
\begin{itemize}
\item علائم حرکتی:
\begin{itemize}
\item ضعف عضلانی: \PNDhl{دیستال بیشتر از پروگزیمال} / \PNDhl{اندام تحتانی بیشتر از فوقانی}.
\item در گیلن باره: درگیری حرکتی بیشتر از حسی.
\item نوروپاتی ناشی از \PNDhl{مواد توکسیک، متابولیک}: درگیری \PNDhl{حسی بیشتر از حرکتی}.
\item \PNDhl{آتروفی عضلات: به‌ویژه در نوروپاتی آکسونال}.
\item کاهش تون و فلاسیدیتی.
\item فاسیکولاسیون، میوکیمی و کرامپ.
\item \PNDhl{نوروپاتی ناشی از اورمی: \PNDen{RLS}}.
\end{itemize}
\item علائم حسی:
\begin{itemize}
\item در پلی‌نوروپاتی‌ها: دستکشی و جورابی.
\item هایپوستزی، آنستزی، \PNDen{Numbness}، پارستزی (حس‌های غیرعادی)، دیس‌استزی (حس‌های ناراحت‌کننده).
\item آتاکسی حسی، تست رومبرگ مثبت، مفاصل نوروپاتیک شارکو.
\end{itemize}
\item تغییرات رفلکسی:
\begin{itemize}
\item \PNDhl{کاهش یا فقدان \PNDen{DTR}}.
\item رفلکس پوستی شکمی و پاسخ کف پایی سالم هستند مگر قوس رفلکسی آسیب دیده باشد.
\end{itemize}
\item در نوروپاتی، فیبرهای نازک درگیر می‌شوند: اختلال رفلکس نداریم چراکه رفلکس مربوط به فیبر ضخیم است.
\end{itemize}
\section*{نوروپاتی عصب مدین}
\begin{itemize}
\item \PNDhl{شایع‌ترین محل آسیب: تونل کارپ} $\longleftarrow$ \PNDhl{\PNDen{CTS}}.
\item \PNDen{RF} برای \PNDen{CTS}: بارداری، \PNDhl{هایپوتیروئیدی}، آکرومگالی، سینوویت، آرتریت روماتوئید، سابقهٔ تروما، فعالیت شغلی، تایپیست.
\item علائم: \PNDhl{علائم حسی زودتر از علائم حرکتی}.
\item درد، هایپوستزی، پارستزی در نواحی توزیع عصب مدین.
\item علائم مثبت بروز می‌کند $\longleftarrow$ از خواب بیدار می‌کند.
\item با تکان دادن دست بهبود می‌یابد.
\item علائم با بلند بردن دست و \PNDunclear{کشش گردنی} بدتر می‌شوند.
\end{itemize}

\PNDpage{۲}
\section*{ادامهٔ نوروپاتی عصب مدین}
\begin{itemize}
\item معاینه:
\begin{itemize}
\item \PNDhl{علامت تینل}: ضربه بر مسیر عصب به مدت ۱ دقیقه.
\item \PNDhl{علامت فالن}: فلکشن مچ دست به مدت ۱ دقیقه.
\item علامت فالن ریورس: اکستنشن مچ دست به مدت ۱ دقیقه.
\item ضعف عضلات مربوطه و آتروفی تنار.
\end{itemize}
\item \PNDhl{تشخیص: \PNDen{EMG}}.
\item درمان: رفع علت زمینه‌ای، اسپلینت مچ دست در شب، تزریق موضعی \PNDen{CS}، جراحی.
\end{itemize}
\section*{نوروپاتی لترال فمورال کوتانیوس}
\begin{itemize}
\item بیماری \PNDhl{مرالژیا پارستتیکا}.
\item فشار بر عصب در ناحیهٔ \PNDhl{لیگامان اینگوئینال}.
\item \PNDen{RF}: \PNDhl{حاملگی، چاقی}.
\item سوزش و پارستزی در سطح \PNDhl{خارجی ـ قدامی} ران.
\item درمان: رفع علت زمینه‌ای، گاباپنتین، جراحی.
\end{itemize}
\section*{نوروپاتی عصب \PNDen{Facial}}
\begin{itemize}
\item ایدیوپاتیک: فلج بل (\PNDen{HSV-1}).
\item ثانویه: ارتشاح، تومور عصب \PNDen{Facial}، تومور \PNDen{CPA}، شکستگی پتروس، تروما به عصب، \PNDhl{گیلن باره}، \PNDhl{لایم}، \PNDhl{سارکوئیدوز}، \PNDhl{ایدز}، \PNDhl{سندرم رامسی هانت}.
\item \PNDen{RF} فلج بل: بارداری، \PNDen{DM}.
\item علائم: درد گوش و اطراف آن، \PNDhl{فلج کامل یا ناقص عضلات یک نیمهٔ صورت}.
\item \PNDhl{اختلال حس چشایی دوسوم قدامی زبان}.
\item اختلال در ترشح بزاق از غدد ساب‌مندیبولر و ساب‌لینگوال در صورت درگیری کوردا تیمپانی.
\item اختلال در ترشح اشک در صورت درگیری عصب.
\item \PNDhl{هایپراکوزیس} در اثر فلج عضلهٔ استاپدیوس $\longleftarrow$ بلند شدن صدا.
\end{itemize}

\PNDpage{۳}
\section*{ادامهٔ نوروپاتی عصب فاسیال}
\begin{itemize}
\item بررسی‌ها: \PNDhl{\PNDen{EMG} بعد از روز ۱۵} $\longleftarrow$ تشخیص، پروگنوز.
\item سیر بیماری: ظرف چند روز تا چند هفته بهبود می‌یابد.
\item اگر هنگام ترمیم عصب، فیبرهای مربوط به بزاق به سمت غدد اشکی بروند:
\item \PNDhl{پدیدهٔ اشک تمساح}: موقع غذا خوردن دچار گریه.
\item درمان: \PNDhl{آسیکلوویر}، \PNDhl{پردنیزولون} (سه روز اول، \PNDen{60 mg/day}، به مدت ۷ روز).
\item عوامل پیش‌آگهی بد:
\begin{itemize}
\item درد شدید گوش در ابتدای بیماری.
\item فلج کامل در ابتدا.
\item یافته‌های \PNDen{EMG}:
\begin{itemize}
\item تغییرات ناشی از آسیب آکسونی (\PNDhl{فیبریلاسیون}).
\item عدم وجود فعالیت الکتریکی ناشی از انقباض عضله.
\item افت پتانسیل حرکتی به کمتر از ۱۰\% فرد سالم.
\end{itemize}
\end{itemize}
\end{itemize}
\section*{پلکسوپاتی براکیال}
\begin{itemize}
\item \PNDen{Thoracic outlet syndrome} (\PNDhl{\PNDen{TOS}}).
\item محل فشار: شبکهٔ براکیال در گردن.
\item دندهٔ گردنی، باندهای فیبروزه، هایپرتروفی عضلهٔ اسکالن قدامی.
\item آسیب اعصاب: \PNDhl{معمولاً ریشهٔ \PNDen{$C_8$} تا \PNDen{$T_1$}}.
\item درد مبهم، ضعف و آتروفی عضلات کوچک دست.
\item آسیب شریان: ایسکمی دست / \PNDhl{تست \PNDen{Adson}}.
\item آسیب ورید: تورم و پرخونی دست.
\item درمان: رفع عامل زمینه‌ای، جراحی.
\end{itemize}
\section*{پلی‌نوروپاتی‌ها}

\PNDpage{۴}
\section*{سندرم گیلن باره}
\begin{itemize}
\item \PNDhl{شایع‌ترین فلج حاد در بالغین جوان} / \PNDhl{\PNDen{$M=F$}}.
\item در فصل سرد شایع‌تر است.
\item تحریک سیستم ایمنی علیه میلین اعصاب محیطی.
\item \PNDhl{\PNDen{anti-GQ1b Ab}} در انواع همراه با \PNDhl{افتالموپلژی}.
\item اتفاقات همراه: عفونت کمپیلوباکتر ژژونی یا مایکوپلاسما پنومونیه، تزریق برخی واکسن‌ها، جراحی.
\item علائم: حسی، حرکتی، احشایی، اتونوم / \PNDhl{حرکتی بیشتر از حسی}.
\item \PNDhl{بیمار تب ندارد} / \PNDhl{علائم قرینه هستند}.
\item علائم حرکتی: \PNDhl{شایع‌ترین شکل آن به صورت بالارونده است}.
\begin{itemize}
\item دیستال اندام تحتانی $\longleftarrow$ پروگزیمال اندام تحتانی $\longleftarrow$ اندام فوقانی $\longleftarrow$ تنه، گردن.
\item \PNDhl{رفلکس مردمک ممکنه از بین می‌رود}.
\item \PNDhl{\PNDen{DTR} عضلات درگیر کم یا از بین می‌روند}.
\item \PNDhl{درگیری عصب \PNDen{Facial} نسبت به بقیهٔ اعصاب کرانیال شایع‌تر است}.
\end{itemize}
\item علائم حسی: کاهش حس، پارستزی انگشتان دست و پا؛ سطح حسی وجود ندارد.
\item در ۵۰\% بیماران \PNDhl{درد عضلانی خصوصاً در کمربند لگنی}.
\item درگیری اتونوم: \PNDhl{شایع} / می‌تواند شدید باشد.
\begin{itemize}
\item تاکی‌کاردی، برادی‌کاردی، آریتمی، فشار خون متغیر.
\item اختلال اسفنکتر (مثانه، روده)، \PNDunclear{اختلال میکسیون}، اختلال عملکرد ریه.
\item ایلئوس پارالیتیک.
\end{itemize}
\item \PNDhl{اختلالات اسفنکتری در آغاز بیماری وجود ندارد} / \PNDhl{افتراق از میلیت عرضی}.
\end{itemize}

\PNDpage{۵}
\section*{ادامهٔ گیلن باره}
\begin{itemize}
\item گیلن باره نوع فیشر: \PNDhl{افتالموپلژی، آتاکسی اندام‌ها، آرفلکسی}.
\item آزمایش‌ها:
\begin{itemize}
\item \PNDen{CSF}: \PNDhl{پس از ۱ هفته} تغییرات دیده می‌شود.
\item \PNDhl{\PNDen{$\uparrow pro$}، بدون \PNDen{$\downarrow$} قند} $\longleftarrow$ \PNDhl{جدایی آلبومینوسیتولوژیک}.
\item کمتر از ۵۰ مونونوکلئر / \PNDhl{\PNDen{PMN} دیده نمی‌شود}.
% A crossed-out Latin abbreviation next to 50 is retained here: PMN.
\end{itemize}
\item الکترومیوگرافی:
\begin{itemize}
\item کاهش شدید سرعت هدایت، بلوک هدایتی، افزایش \PNDen{Latency}.
\item از بین رفتن موج \PNDen{F} و \PNDen{H}.
\item اگر آسیب به آکسون: فیبریلاسیون.
\item یافته‌های \PNDen{EMG} دیروقت و غیر اختصاصی هستند.
\end{itemize}
\item در صورت شک به درگیری اتونوم: \PNDen{EKG}، اسپیرومتری در صورت وجود علائم تنفسی / \PNDen{ABG} جهت بررسی عملکرد ریه.
\item \PNDhl{تشخیص گیلن باره در مراحل اولیه بالینی است}.
\item سیر بیماری: \PNDhl{طی ۲ تا ۴ هفته پیشرفت} $\longleftarrow$ سپس ثابت می‌گردد $\longleftarrow$ در نهایت بهبودی یابد.
\item \PNDhl{شک به گیلن باره $\longleftarrow$ بستری و انجام بررسی‌های بیشتر}.
\item درمان:
\begin{itemize}
\item \PNDhl{پلاسمافرز}: \PNDen{200--250 cc/kg} طی ۴–۶ جلسه تقسیم شود.
\item \PNDhl{\PNDen{IVIG}}: \PNDen{400 mg/kg} به مدت ۵ روز.
\item بستری \PNDen{ICU} در صورت نیاز.
\item اندیکاسیون‌ها انتوباسیون: اختلال بلع، ضعف عضلات تنفسی، کاهش \PNDhl{ظرفیت حیاتی (\PNDen{VC}) به کمتر از \PNDen{10 ml/kg}}.
\item درمان اختلالات اتونوم:
\begin{itemize}
\item \PNDen{$BP\downarrow$}: مایع‌درمانی + داروهای بالابرندهٔ فشار.
\item \PNDen{$BP\uparrow$}: \PNDhl{لابتالول \PNDen{IV}}.
\item آریتمی: داروی ضد آریتمی.
\end{itemize}
\end{itemize}
\end{itemize}

\PNDpage{۶}
\section*{ادامهٔ گیلن باره}
\begin{itemize}
\item هایپوناترمی: ابتدا محدودیت مایعات $\longleftarrow$ سپس محلول سدیم.
\item در فلج شدید برای جلوگیری از \PNDen{DVT}: \PNDhl{هپارین}.
\item فیزیوتراپی تنفسی و اندام‌ها.
\item \PNDhl{\PNDen{CS} در گیلن باره جایی ندارد}.
\item پروگنوز: ۷۵\% خوب می‌شوند، ۲۰\% عارضه‌دار می‌شوند، ۵\% مرگ.
\item \PNDhl{پیش‌آگهی بد: سن بالا، درگیری آکسونال، نیاز به انتوباسیون، پیشرفت سریع علائم}.
\item علل مرگ: نارسایی تنفسی، آسپیراسیون، آمبولی، زخم بستر، عفونت، آریتمی، \PNDen{$BP\downarrow\uparrow$}.
\item درگیری ایزولهٔ حسی $\longleftarrow$ \PNDen{R/O} گیلن باره.
\end{itemize}
\section*{پلی‌نوروپاتی دمیلینیزان التهابی مزمن ـ \PNDen{CIDP}}
\begin{itemize}
\item آسیب اولیه در میلین است / سیر بیماری آهسته‌تر از گیلن باره.
\item بیماری \PNDhl{حداقل بیشتر از ۶ هفته} و \PNDhl{اغلب بیشتر از ۱۲ هفته} قبل از تشخیص وجود داشته است.
\item انواع بالینی:
\begin{itemize}
\item پیشرفت تدریجی و مزمن.
\item بدتر شدن بیمار به صورت پله‌ای.
\item عودکننده: به صورت دوره‌های بهبودی و تشدید.
\end{itemize}
\item علائم:
\begin{itemize}
\item \PNDhl{ممکن است برخلاف گیلن باره درگیری اندام فوقانی بیش از تحتانی باشد}.
\item ممکن است علائم بیماری \PNDhl{غیرقرینه} باشد.
\item ممکن است \PNDhl{اعصاب ضخیم و قابل لمس} باشد.
\end{itemize}
\item \PNDhl{تغییرات \PNDen{CSF} و \PNDen{EMG} شبیه گیلن باره} است.
\item درمان:
\begin{itemize}
\item \PNDhl{پردنیزولون}: \PNDen{40--100 mg daily/high dose}.
\item \PNDhl{\PNDen{IVIG}}.
\item \PNDhl{پلاسمافرز}.
\item آزاتیوپرین، سیکلوفسفامید $\longleftarrow$ در صورت عدم پاسخ.
\end{itemize}
\end{itemize}

\PNDpage{۷}
\section*{\PNDen{DDX} بیماری \PNDen{CIDP} از گیلن باره}
\begin{itemize}
\item اغلب \PNDen{CIDP} بیشتر از ۱۲ هفته طول می‌کشد.
\item ممکن است درگیری از اندام فوقانی آغاز شود.
\end{itemize}
\section*{پلی‌نوروپاتی‌ها: اندوکراین}
\begin{itemize}
\item \PNDhl{\PNDen{DM} (مهم‌ترین)}، اختلالات تیروئید، آکرومگالی.
\item در حاشیه: \PNDen{CTS} (با شاخه‌هایی به موارد بالا).
\end{itemize}
\section*{پلی‌نوروپاتی دیفتریایی}
\begin{itemize}
\item به دلیل نوروتوکسین کورینه‌باکتریوم دیفتریه $\longleftarrow$ آسیب میلین اعصاب محیطی.
\item علائم:
\begin{itemize}
\item نوع کرانیال: یک ماه بعد از عفونت اولیه شروع می‌شود.
\item تاری دید، فلج عضلهٔ حلق، چشم، کام نرم، صورت، حلق و دیافراگم.
\item اختلال رفلکس مردمک در نگاه به نزدیک، فلج تنفسی.
\item نوع اندام‌ها: ۲–۳ ماه بعد از عفونت اولیه.
\item شبیه گیلن باره: پلی‌نوروپاتی حسی و حرکتی دیستال و قرینه.
\item ممکن است فلج غیرقرینه باشد و درگیری \PNDen{prox} بیشتر از \PNDen{distal}.
\end{itemize}
\item درمان: \PNDen{AB}، آنتی‌توکسین، حمایتی.
\end{itemize}
\section*{پلی‌نوروپاتی ناشی از جذام}
\begin{itemize}
\item لپر توبروکلوئید: ماکول یا پاپول هایپوپیگمانته که حس درد و حرارت در آن‌ها مختل است.
\item تعریق در این بیماران وجود ندارد.
\item اختلال حسی: در اعصاب دیجیتال، سورال، رادیال، خلف گوشی.
\item اختلال حرکتی در عضلات اعصاب پرونئال و اولنار.
\item اعصاب درگیر در لمس بزرگ و ضخیم هستند.
\end{itemize}

\PNDpage{۸}
\section*{ادامهٔ پلی‌نوروپاتی ناشی از جذام}
\begin{itemize}
\item لپر لپروماتوز: درگیری منتشر حسی است / حس درد و حرارت.
\item زخم‌های تروفیک، سیانوز، عدم تعریق دست و پا.
\item \PNDen{DTR} سالم است مگر در موارد شدید.
\item مناطق غیرپوشیدهٔ بدن بیشتر درگیر می‌شود: بینی، گوش، گونه، پشت دست‌ها، ساعد، پا، سطح خارجی ساق پا.
\end{itemize}
\section*{پلی‌نوروپاتی هرپس زوستر}
\begin{itemize}
\item فعال شدن مجدد ویروس در گانگلیون ریشهٔ خلفی.
\item \PNDen{RF}: سن بالا، نقص ایمنی مثل لوکمی و لنفوم.
\item علائم: درد در درماتوم درگیر / سپس در ۲–۵ روز راش وزیکولر قرمز رنگ با روند راش درد کم می‌شود.
\item درد خنجری است.
\item نورالژی بعد از هرپس ممکن است رخ دهد.
\item مناطق درگیر:
\begin{itemize}
\item توراسیک: شایع‌ترین منطقه.
\item شاخهٔ اول عصب تری‌ژمینال.
\item عصب \PNDen{Facial} (سندرم رامسی هانت).
\end{itemize}
\item \PNDen{CSF}: \PNDen{$\uparrow$} سلول و \PNDen{pro}.
\item درمان: مسکن، \PNDen{CS} (درمان کاهش‌دهندهٔ نورالژی نیست)، آسیکلوویر (پیشگیری از نورالژی).
\item داروهای مؤثر در تسکین درد نورالژی بعد از هرپس:
\begin{itemize}
\item \PNDen{TCA}: آمی و نورتریپتیلین.
\item ضدتشنج: گاباپنتین، کاربامازپین.
\item کپسایسین موضعی.
\end{itemize}
\end{itemize}

\PNDpage{۹}
\section*{مسمومیت با سرب}
\begin{itemize}
\item باتری‌سازی، کشتی‌سازی، رنگ‌سازی.
\item علائم: کودکان: درگیری \PNDen{CNS} / \PNDhl{بزرگسالان: درگیری \PNDen{PNS}}.
\item علائم درگیری \PNDen{PNS}: \PNDhl{اندام فوقانی بیشتر از اندام تحتانی}.
\item \PNDhl{بیشتر عصب رادیال} را درگیر می‌کند: \PNDhl{افتادگی مچ دست}.
\item \PNDhl{ضعف عضلانی غیرقرینه}.
\item درگیری غیرنورولوژیک: آنمی، یبوست، دردهای کولیکی شکمی، \PNDhl{تغییر رنگ لثه}، درگیری کلیه.
\item درمان: \PNDhl{\PNDen{EDTA}}، \PNDhl{\PNDen{BAL}}، \PNDhl{پنی‌سیلامین}.
\end{itemize}
\section*{مسمومیت با تالیوم}
\begin{itemize}
\item \PNDhl{ابتدا با پلی‌نوروپاتی حسی‌حرکتی سریعاً پیش‌رونده شبیه به سندرم گیلن باره}.
\item \PNDhl{علائم گوارشی قبل، همزمان یا بعد از علائم نورولوژیک}.
\item \PNDhl{علائم پوستی: ریزش مو، پوست لکه‌لکه}.
\item تشخیص: شناسایی \PNDhl{تالیوم در ادرار} و بافت‌های بدن / \PNDen{CSF} مانند گیلن باره است.
\item درمان: حمایتی.
\end{itemize}
\section*{پورفیری}
\begin{itemize}
\item اختلال در بیوسنتز هم.
\item عوامل شعله‌ساز:
\begin{itemize}
\item داروها (باربیتورات‌ها، کلردیازپوکساید، سولفونامیدها، \PNDen{OCP}، متیل‌دوپا، ایمی‌پرامین، کلروکین، ریفامپین، الکل).
\item گرسنگی، بارداری، استرس، عفونت.
\end{itemize}
\item علائم شکمی: ابتدا \PNDhl{درد کولیکی شکم، یبوست، تب، استفراغ، لکوسیتوز، بی‌قراری و بی‌اشتهایی}.
\item علائم نورولوژیک: \PNDhl{نوروپاتی شبیه گیلن باره} می‌تواند \PNDhl{دردناک، قرینه یا غیرقرینه} باشد.
\item ممکن است درگیری فوقانی بیشتر از تحتانی و \PNDen{prox} بیشتر از دیستال باشد.
\item علائم اتونوم: \PNDhl{تاکی‌کاردی}؛ ممکن است \PNDhl{\PNDunclear{منجر به فلج روده گردد}}.
\item درمان: \PNDhl{گلوکز و هماتین}.
\end{itemize}
\endgroup
\ifPNDStandalone
  \expandafter\enddocument
\fi
